\section{Models and Evaluation}\label{sec:exploit}

This section sets out a simple framework for the speed of price discovery, the common trading
and evaluation environment, and the four forecasting models. Throughout, $T$ denotes the release
minute, $P_\tau$ the close of the one-minute bar labelled $\tau$, $p_\tau=\ln P_\tau$, and
$\mathcal{F}_{T}$ the information available at the close of bar $T$, one minute after the
release. All returns are in basis points (bps).

\subsection{A partial-adjustment framework}\label{sec:framework}

Let $v_T$ denote the change in the fundamental value of the asset implied by the news in a
release, and suppose that it is proportional to the family signal in equation~\eqref{eq:x},
$v_T=\gamma\,x_T$, where $\gamma$ is the full-information price impact of a one-standard-deviation
surprise. Prices adjust toward the new value gradually. We write the cumulative reaction after
$h$ minutes as
\begin{equation}\label{eq:pa}
  R_{T,h} = p_{T+h-1}-p_{T-1} = \Lambda_h\,\gamma\,x_T + u_{T,h},
  \qquad 0\le \Lambda_1 \le \Lambda_2 \le \dots \le 1,
\end{equation}
where $\Lambda_h$ is the share of the news priced after $h$ minutes and $u_{T,h}$ is noise
unrelated to the surprise, $\mathbb{E}[u_{T,h}\mid x_T]=0$. Under efficient and instantaneous
price discovery, $\Lambda_1=1$. Using the decomposition in equation~\eqref{eq:decomp}, the
tradeable drift from the first post-release price is
\begin{equation}\label{eq:drift_model}
  D_{T,h} = R_{T,h}-R_{T,1} = (\Lambda_h-\Lambda_1)\,\gamma\,x_T + (u_{T,h}-u_{T,1}).
\end{equation}
Equation~\eqref{eq:drift_model} delivers the two quantities this paper measures. The first is the
share of the news that is \emph{not} priced within the first minute,
\begin{equation}\label{eq:theta}
  \theta_h \equiv 1-\frac{\Lambda_1}{\Lambda_h}
  = \frac{\beta^{D}_h}{\beta^{R}_h},
  \qquad
  \beta^{R}_h=\frac{\operatorname{Cov}(R_{T,h},x_T)}{\operatorname{Var}(x_T)},\quad
  \beta^{D}_h=\frac{\operatorname{Cov}(D_{T,h},x_T)}{\operatorname{Var}(x_T)},
\end{equation}
which is identified from two predictive regressions. The second is the predictability of the
drift from public information: since $x_T\in\mathcal{F}_T$,
\begin{equation}\label{eq:predict}
  \mathbb{E}\left[D_{T,h}\mid\mathcal{F}_T\right] = \beta^{D}_h\,x_T,
\end{equation}
so a strategy that takes a position in the direction of $x_T$ at the close of bar $T$ earns
$\beta^{D}_h\,\mathbb{E}|x_T|$ per trade before costs. The drift is exploitable only if this
exceeds the round-trip cost $c$,
\begin{equation}\label{eq:exploit}
  \beta^{D}_h\,\mathbb{E}\left|x_T\right| > c .
\end{equation}
Condition~\eqref{eq:exploit} separates two reasons why a predictable drift may not be traded: the
news may be priced almost entirely within a minute ($\theta_h\approx 0$), or the remaining drift
may be small relative to the minimum price increment, which sets a floor on $c$. The average CPI
paths in Section~\ref{sec:discovery} imply $\theta_{30}\approx 0.3$ in \ES, \NQ\ and \YM\ and
$\theta_{30}\approx 0.2$ in \ZN; the estimated drift slopes of about 5.9, 8.5, 4.4 and 0.8 bps per
standard deviation, against round-trip costs of about 1.7, 0.7, 1.0 and 3.0 bps, satisfy
condition~\eqref{eq:exploit} in the three equity markets and violate it in \ZN.

\subsection{Trading environment}\label{sec:environment}

\paragraph{Timing and returns.}
Every model forms its forecast with information in $\mathcal{F}_T$ and opens a position
$w_T\in[-2,2]$ (in contracts) at the close of bar $T$ (Figure~\ref{fig:timing}). A trade on
window $k$ opens at minute $T+a_k$ and closes 30 minutes later, with gross return
\begin{equation}\label{eq:trade}
  r^{g}_{T,k} = w_{T,k}\,\left(p_{T+a_k+29}-p_{T+a_k-1}\right)\times 10^4 ,
\end{equation}
with $a_1=1$ for the single-window models, so that $r^{g}_{T,1}=w_{T,1}D_{T,30}$. The last
pre-release price $p_{T-1}$ is never an entry price.

\begin{figure}[htbp]
  \centering
  \includegraphics[width=0.85\linewidth]{diagram_timing}
  \caption{Timing of a trade. The release occurs at $T$; the close of bar $T-1$ is the last
  pre-release price; the jump is priced within bar $T$; the position is opened at the close of
  bar $T$ and closed at the close of bar $T+29$.}
  \label{fig:timing}
\end{figure}

\paragraph{Costs.}
For a contract with tick size $\delta$ (in price points) and point value $M$ (in dollars), trading
$n$ ticks of slippage per round trip plus a commission $\kappa$ costs, in basis points of the
entry price $P$,
\begin{equation}\label{eq:cost}
  c(P) = \frac{n\,\delta + \kappa/M}{P}\times 10^4 ,
\end{equation}
and the net return is $r_{T,k}=r^{g}_{T,k}-|w_{T,k}|\,c(P_{T+a_k-1})$. We report $n\in\{0,
\tfrac12, 1\}$ with $\kappa=0$ (gross, a quarter and a half tick per side) and $n=2$, $\kappa=\$4.50$
(the base case).

\paragraph{Performance.}
Let $d(t)$ be the trading session of time $t$ (18:00--17:00 Eastern Time, dated by the close). The
daily P\&L of a book is $\pi_s=\sum_{(T,k):\,d(T+a_k)=s} r_{T,k}$, zero on sessions without a trade.
Over $S$ sessions,
\begin{equation}\label{eq:sharpe}
  \widehat{SR} = \sqrt{252}\,\frac{\bar\pi}{\hat\sigma_\pi},\qquad
  \mathrm{MDD} = \min_{s}\Big(\sum_{j\le s}\pi_j-\max_{u\le s}\sum_{j\le u}\pi_j\Big),\qquad
  \hat\beta = \frac{\widehat{\operatorname{Cov}}(\pi_s,\pi^{BH}_s)}{\widehat{\operatorname{Var}}(\pi^{BH}_s)},
\end{equation}
where $\pi^{BH}_s$ is the daily return of holding one contract, with roll jumps set to zero. The
single-model results in Section~\ref{sec:results}.A report monthly Sharpe ratios, with
$\sqrt{12}$ replacing $\sqrt{252}$.

\paragraph{Out-of-sample estimation.}
Every estimated quantity---line signs, the release universe, regression slopes, gate thresholds
and position-size references---is re-estimated at the start of each year $y$ using only releases
before January 1 of year $y$. We write $\mathcal{H}_y$ for this history. The single-model results
(Section~\ref{sec:results}.A) use history from 2016 and evaluate 2019--2026 at the base cost. The
comparison of the four models (Sections~\ref{sec:results}.C--F) uses the split in
Table~\ref{tab:split} and reports the validation years (2018--2020) and the test years (2021--2026)
separately.

\subsection{Model 1: Consensus surprises}\label{sec:model_surprise}

\paragraph{Line signs.} For each market and line $l$, the sign $s_{l,y}$ in
equation~\eqref{eq:x} is the sign of the OLS slope of the first-minute jump on the standardized
surprise over $\mathcal{H}_y$,
\begin{equation}\label{eq:sign}
  R_{T,1} = a_l + b_l\, z_{l,T} + e_{T},\qquad s_{l,y}=\operatorname{sign}\big(\hat b_{l}\big),
\end{equation}
estimated when at least 24 releases are available.

\paragraph{Release universe.} For each family $f$ and year $y$, an impact screen compares the
family's 30-minute reaction to that of other families released at the same clock time,
\begin{equation}\label{eq:impact}
  I_{f,y} = \frac{\operatorname{mean}_{T\in\mathcal{T}_{f}\cap\mathcal{H}_y}\,|R_{T,30}|}
  {\operatorname{mean}_{T\in\mathcal{C}_{f}\cap\mathcal{H}_y}\,|R_{T,30}|},
\end{equation}
where $\mathcal{T}_f$ are the family's release times and $\mathcal{C}_f$ the release times of other
families at the family's modal clock time on which $f$ is not released. The ten families with the
largest $I_{f,y}$ (at least 30 releases each), plus CPI, form the universe $\mathcal{U}_y$.

\paragraph{Naive book.} For every $f\in\mathcal{U}_y$ released at $T$,
\begin{equation}\label{eq:naive}
  w_{f,T} = \operatorname{sign}\left(x_{f,T}\right),\qquad
  w_T = \frac{1}{|\mathcal{U}_T|}\sum_{f\in\mathcal{U}_T} w_{f,T},
\end{equation}
where $\mathcal{U}_T$ is the set of universe families released at the same minute, so no minute is
traded twice.

\paragraph{CPI model.} For CPI, the drift slope in equation~\eqref{eq:predict} is estimated on
the earlier CPI releases $\mathcal{H}_T=\{t<T\}$,
\begin{equation}\label{eq:cpimodel}
  \hat\beta^{D}_{T} = \arg\min_{b}\sum_{t\in\mathcal{H}_T}\big(D_{t,30}-b\,x_{t}\big)^2,
  \qquad \hat r_T=\hat\beta^{D}_{T}\,x_T,
\end{equation}
requiring at least 24 earlier releases, and the position implements
condition~\eqref{eq:exploit} release by release:
\begin{equation}\label{eq:cpitrade}
  w_T = \operatorname{sign}(\hat r_T)\,\mathbf{1}\big\{|\hat r_T| > c(P_{T})\big\}.
\end{equation}

\subsection{Model 2: Price-path analogues (DTW ladder)}\label{sec:model_dtw}

\paragraph{Paths.} For a trade entering at minute $e=T+a_k$, the query path is the log price over
the $L+1=31$ minutes ending at the entry, standardized so that only its shape matters,
\begin{equation}\label{eq:path}
  q_j = \frac{p_{e-L+j-1}-\bar p}{\hat\sigma_p},\qquad j=1,\dots,L+1 ,
\end{equation}
where $\bar p$ and $\hat\sigma_p$ are the mean and standard deviation of the $L+1$ log prices.

\paragraph{Distance.} For two standardized paths $q$ and $c$ of length $n$, the dynamic time
warping distance is computed by the recursion
\begin{equation}\label{eq:dtw}
  g(i,j) = (q_i-c_j)^2 + \min\{g(i-1,j),\,g(i,j-1),\,g(i-1,j-1)\},
  \qquad |i-j|\le b,
\end{equation}
with $g(0,0)=0$, $g(i,0)=g(0,j)=\infty$ and $d(q,c)=g(n,n)$. The Sakoe--Chiba band $b=5$ minutes
allows the reaction to unfold up to five minutes faster or slower than in the matched past
episode \citep{Sakoe1978, Berndt1994}.

\paragraph{Neighbours and forecasts.} The candidate set $\mathcal{N}_e$ contains the paths of
earlier releases at the same clock time and window $k$, within a three-year lookback, whose own
trade closed before $e$, limited to the 600 most recent. With $\mathcal{K}_e\subset\mathcal{N}_e$
the $K=15$ candidates with the smallest distances $d_j$, ages $\tau_j$ in years and realized
next-30-minute returns $y_j$, the weights and the two forecasts are
\begin{equation}\label{eq:dtwforecast}
  \omega_j = \frac{d_j^{-1}\,2^{-\tau_j/H}}{\sum_{i\in\mathcal{K}_e} d_i^{-1}\,2^{-\tau_i/H}},\qquad
  \mu_e = \sum_{j\in\mathcal{K}_e}\omega_j\,y_j,\qquad
  \sigma_e = \Big(\sum_{j\in\mathcal{K}_e}\omega_j\,(y_j-\mu_e)^2\Big)^{1/2},
\end{equation}
with a recency half-life of $H=2$ years. The direction forecast $\mu_e$ is the analogue estimate of
$\mathbb{E}[D\mid\text{path}]$ and the dispersion forecast $\sigma_e$ measures how much the analogues
disagree.

\paragraph{Gate and size.} With $Q_\alpha$ the $\alpha$-quantile of $\mu$ over the fitting years and
$\bar\sigma$ the median of $\sigma$,
\begin{equation}\label{eq:dtwtrade}
  w_e = \operatorname{sign}(\mu_e)\,\mathbf{1}\big\{\mu_e\le Q_{0.3}\ \text{or}\ \mu_e\ge Q_{0.7}\big\}
  \cdot \min\Big\{\frac{\bar\sigma}{\sigma_e},\,2\Big\}.
\end{equation}
The book trades five back-to-back windows after each release of eight families, $a_k\in\{1,31,61,
91,121\}$, with the forecasts in equation~\eqref{eq:dtwforecast} rebuilt for each window.

\subsection{Model 3: Pattern-based event games}\label{sec:model_games}

The event-game model applies the analogue logic to longer paths around the release. For each
release $T$ and anchor $T_0=T+o$, $o\in\{-60,-30,0,30,60\}$ minutes, the path is the log return
over $[T_0,T_0+60]$ sampled every $\Delta=3$ minutes and standardized as in
equation~\eqref{eq:path}, and the target is the next 30-minute return,
\begin{equation}\label{eq:game}
  y_{T,o} = p_{T_0+90}-p_{T_0+60}.
\end{equation}
Neighbours come from a pool $\mathcal{P}$ of release types (Jobless Claims, JOLTS and PCE) at the
same clock time and anchor, within a two-year lookback, with a DTW band of five sampling steps. The
factor is the distance- and recency-weighted average of the neighbours' targets, as in
equation~\eqref{eq:dtwforecast}, and the book trades
\begin{equation}\label{eq:gametrade}
  w_{T,o}=\operatorname{sign}\big(\mu_{T,o}\big)
\end{equation}
one contract, entering at $T_0+60$. The pool, sampling frequency and lookback were chosen by a grid
search over 18 combinations on the full sample.

\subsection{Model 4: Machine-learning classifier}\label{sec:model_ml}

The classifier forecasts the direction of the 30-minute return from a feature vector $\phi_T$
that combines surprise, price-path and momentum information. With $\hat F_y$ trained on
$\mathcal{H}_y$,
\begin{equation}\label{eq:ml}
  \hat\rho_T = \Pr\nolimits_{\hat F_y}\left(D_{T,30}>0 \mid \phi_T\right),\qquad
  w_T = g\big(\hat\rho_T\big)\in[-1,1],
\end{equation}
where $g$ maps the predicted probability into a signed position that is larger when the model is
more confident; we keep the model's positions exactly as produced. The \emph{executable} version
uses the trade return in equation~\eqref{eq:trade} with $a_1=1$. The \emph{as-reported} version
enters at the last pre-release price,
\begin{equation}\label{eq:asreported}
  r^{g,\text{rep}}_{T} = w_T\,\big(p_{T+28}-p_{T-1}\big)\times 10^4 = w_T\,\big(R_{T,1}+D_{T,29}\big)\times 10^4 ,
\end{equation}
so the difference between the two isolates the value of the untradeable first-minute jump
$R_{T,1}$, which by equation~\eqref{eq:pa} contains the share $\Lambda_1$ of the news.

\subsection{Inference}\label{sec:inference}

\paragraph{Random-sign test.} Holding the trades, sizes and costs of a book fixed, we draw
$\epsilon_{T,k}\in\{-1,+1\}$ with equal probability and form
$\tilde r_{T,k}=\epsilon_{T,k}\,|w_{T,k}|\,(p_{T+a_k+29}-p_{T+a_k-1})\times10^4-|w_{T,k}|\,c$. With
$\widehat{SR}^{(b)}$ the Sharpe ratio of the $b$-th of $B=1{,}000$ draws,
\begin{equation}\label{eq:randsign}
  p_{\text{RS}} = \frac{1}{B}\sum_{b=1}^{B}\mathbf{1}\big\{\widehat{SR}^{(b)}\ge\widehat{SR}\big\},
\end{equation}
which tests whether the direction adds value beyond being in the market at release times.

\paragraph{Block bootstrap.} Confidence intervals for $\widehat{SR}$ resample blocks of three
consecutive months of returns with replacement \citep{Kunsch1989}, which preserves serial
dependence in the monthly returns.

\paragraph{Leak test.} For the analogue models we deliberately extend the query path in
equation~\eqref{eq:path} five minutes past the entry. Because the extended path overlaps the
target, performance must rise sharply if the timing is correct; an unchanged result would indicate
that the real book already uses future prices.

\subsection{Combining the models}\label{sec:model_portfolio}

Let $\pi_{i,s}$ be the daily net P\&L of book $i$ at a quarter tick per side and $\hat\sigma_i$ its
standard deviation over the validation years 2018--2020. An equal-risk portfolio of a set of books
$\mathcal{B}$ has daily return
\begin{equation}\label{eq:portfolio}
  \pi^{P}_s = \sum_{i\in\mathcal{B}} \lambda_i\,\pi_{i,s},\qquad
  \lambda_i=\frac{\hat\sigma_i^{-1}}{\sum_{j\in\mathcal{B}}\hat\sigma_j^{-1}},
\end{equation}
and is scored on the test years only. For uncorrelated books with equal Sharpe ratios $SR$, the
portfolio Sharpe ratio is $\sqrt{|\mathcal{B}|}\,SR$; the gain from combining the four models
therefore depends on how correlated their daily returns are and on how much each loses after costs.
